\section{Results}
\label{sec:results}

\subsection{Verification and validation}
\label{sec:verification}

\begin{figure}[htbp]
\centering
\includegraphics[width=\textwidth]{fig03_validation_xf1989}
\caption{Verification against the digitised industrial case of Xu and Froment \cite{xu1989b} (symbols) with the digitised
outer-wall temperature as boundary condition: (a) methane and CO$_2$ conversion, (b) process-gas and inner-wall temperatures,
(c) pressure and (d) temperature residuals, for the matched film coefficient $\alpha_i = \alphaIxf{}$ W m$^{-2}$ K$^{-1}$
(solid lines) and for the Leva--Grummer correlation with $f_{htg} = 1$ (dotted). Source: \texttt{data/literature\_validation/digitized/},
\texttt{xf1989\_fit.yaml}, \rdtversion.}
\label{fig:xf}
\end{figure}

\begin{table}[htbp]
\centering
\caption{Verification of the tube model against the digitised Fig.~3 of Xu and Froment \cite{xu1989b}: RMSE of methane conversion $x_{\mathrm{CH_4}}$, of the conversion increment $x(z)-x(0.5\,\mathrm{m})$, of the process-gas and inner-wall temperatures ($z\ge0.3$~m) and of pressure. Source: \texttt{data/literature\_validation/xf1989\_fit.yaml} (verification\_matched\_alpha).}
\label{tab:xf}
\small
\resizebox{\textwidth}{!}{%
\begin{tabular}{lrrrrr}
\toprule
Configuration & RMSE $x_{\mathrm{CH_4}}$ & RMSE $\Delta x$ & RMSE $T_g$ (K) & RMSE $T_{wi}$ (K) & RMSE $p$ (bar) \\
\midrule
Leva--Grummer, $f_{htg}=1$, $\eta=0.1$ & 0.160 & nan & 75.4 & 37.5 & 0.04 \\
matched $\alpha_i$ = 348 W m$^{-2}$ K$^{-1}$, $\eta=0.1$ & 0.037 & 0.020 & 13.7 & 16.7 & 0.25 \\
matched $\alpha_i$, $\eta$ = 0.05 / 0.1 (adopted) & 0.039 & 0.015 & 14.0 & 16.7 & 0.25 \\
matched $\alpha_i(z)$ profile, $\eta$ = 0.05 / 0.1 & 0.046 & 0.017 & 17.6 & 16.9 & 0.26 \\
\bottomrule
\end{tabular}}
\par\smallskip\footnotesize Back-calculated $\alpha_i$: median 348, interquartile range 313--363 W m$^{-2}$ K$^{-1}$ over 0.5--11 m.
\end{table}


With the digitised outer-wall temperature prescribed, the tube model's response is controlled almost entirely by the bed-side film
coefficient, and the first finding of the study is that the heat-transfer chain printed in Part II of Xu and Froment
\cite{xu1989b} does not reproduce their own figure. Evaluated as printed, the De~Wasch--Froment wall coefficient with the
Yagi--Kunii dynamic terms gives $\alpha_i \approx$ 3100--3700 W m$^{-2}$ K$^{-1}$ at the particle Reynolds number of the case, and
the Leva--Grummer correlation gives 940--1580 W m$^{-2}$ K$^{-1}$; the value back-calculated from the digitised wall and gas
temperatures and the wall conduction is \alphaIxf{} W m$^{-2}$ K$^{-1}$ (interquartile range 313--363). Used as printed, the
correlations overpredict the gas temperature by 70--115 K and the conversion by 0.14--0.22 (Table~\ref{tab:xf}, first row). Both
dynamic terms are linear in $Re\,Pr$ and were fitted at Reynolds numbers one order of magnitude lower than those of an industrial
tube, so the extrapolation drives the effective radial conductivity to values that no measurement supports. A displaced decimal
in the printed coefficient of the effective conductivity would resolve the discrepancy, but we treat this as a hypothesis
\todo{check the coefficients against De Wasch and Froment (1972) and Froment and Bischoff before asserting a misprint}.

The Xu--Froment case is therefore used as a \emph{verification} of the tube model with a matched film coefficient rather than as a
validation of the heat-transfer correlation. With $\alpha_i = \alphaIxf{}$ W m$^{-2}$ K$^{-1}$ and Latham's inlet effectiveness
profile, the model reproduces the digitised profiles with an RMSE of \xfRMSET{} K in gas temperature, 16.7 K in inner-wall
temperature and \xfRMSEinc{} in the conversion increment along the tube (Fig.~\ref{fig:xf}, Table~\ref{tab:xf}); the residual
offset in absolute conversion of about 0.04 originates in the first half metre, where the digitised curve and the inlet rule
differ \todo{comment on the inlet-region offset}. This verifies kinetics, thermodynamics, the energy balance and the wall
conduction chain. The pressure RMSE of \SI{0.25}{bar} with the matched coefficient, against \SI{0.04}{bar} with the Leva--Grummer run,
reflects the different gas temperature and hence density along the tube rather than a change in the friction model
\todo{confirm this explanation of the pressure residual}.

\begin{figure}[htbp]
\centering
\includegraphics[width=\textwidth]{fig05_latham_calibrated}
\caption{Coupled furnace--tube model after calibration ($\alpha_{top}$ fixed) on the Latham plant cases: axial profiles of
furnace-gas, outer-wall, inner-wall and process-gas temperatures for Plants A, B and C2 with the measured wall temperatures at the
two peep-hole elevations, the outlet temperature and the flue-gas outlet temperature (error bars $\pm 1.96\sigma$), and the parity of
predicted against measured temperatures (dotted lines $\pm 10$ K). Source: \texttt{latham2008\_plant\_cases.csv},
\texttt{latham\_fit.yaml}, \rdtversion.}
\label{fig:latham}
\end{figure}

\begin{table}[htbp]
\centering
\caption{Latham plant cases (thesis Appendix H, absolute temperatures) and calibration results: measured inlet and outlet conditions, model errors (predicted minus measured) of the all-cases fit with $\alpha_{top}$ fixed, and leave-one-case-out held-out errors. Plant C1 is excluded from the fit (inconsistent feed flow). Source: \texttt{latham2008\_plant\_cases.csv}, \texttt{latham\_fit.yaml} (calibration\_alpha\_top\_fixed).}
\label{tab:latham}
\small
\resizebox{\textwidth}{!}{%
\begin{tabular}{lrrrrrrrrrrrrr}
\toprule
Case & Rate (\%) & $T_{in}$ (K) & $p_{in}$ (bar) & Feed (kmol h$^{-1}$) & S/C & $T_{out}$ (K) & $T_{w,u}$ (K) & $T_{w,l}$ (K) & $\Delta T_{out}$ & $\Delta T_{w,u}$ & $\Delta T_{w,l}$ & CV $\Delta T_{w,u}$ & CV $\Delta T_{w,l}$ \\
\midrule
Plant A & 99 & 884.5 & 30.06 & 23.5 & 2.89 & 1105.5 & 1095.9 & 1132.0 & +3.6 & -8.4 & +0.2 & -12.8 & +0.7 \\
Plant B & 92 & 887.0 & 29.44 & 21.1 & 2.79 & 1107.2 & 1082.3 & 1130.8 & +2.2 & -3.6 & -1.1 & -5.4 & -1.7 \\
Plant C1 & 96 & 890.9 & 30.27 & 28.7 & 2.85 & 1122.0 & 1111.2 & 1151.5 & -- & -- & -- & -- & -- \\
Plant C2 & 71 & 879.9 & 28.09 & 15.9 & 2.86 & 1108.9 & 1075.7 & 1126.2 & +3.9 & +12.2 & -0.5 & +18.2 & -0.1 \\
\bottomrule
\end{tabular}}
\par\smallskip\footnotesize Per-tube feed; S/C on total hydrocarbon carbon; upper/lower wall temperatures at 3.66 and 8.53 m of 12.5 m; thesis uncertainties: 2~K outlet, 3~K wall, 8~K flue gas.
\end{table}

\begin{table}[htbp]
\centering
\caption{Calibrated furnace and heat-transfer parameters on Plants A, B and C2 with 95\,\% $t$-intervals from the Jacobian covariance. The four-parameter fit has $\alpha_{top}$ at its bound with $r(L_q,\alpha_{top})=0.985$; the adopted fit fixes $\alpha_{top}=0.182$. Source: \texttt{latham\_fit.yaml}.}
\label{tab:cal}
\small
\begin{tabular}{llrlr}
\toprule
Fit & Parameter & Value & 95\,\% interval & s.e. \\
\midrule
four parameters & ${F_{gt}}$ & 0.316 & [0.286, 0.346] & 0.014 \\
four parameters & ${L_q}$ & 0.497 & [0.329, 0.665] & 0.081 \\
four parameters & ${\alpha_{top}}$ & 0.4 & [-0.261, 1.061] & 0.32 \\
four parameters & ${f_{htg}}$ & 1.74 & [1.427, 2.051] & 0.15 \\
four parameters & $\chi^2$ / dof & 30.5 & 23 & reduced 1.32 \\
$\alpha_{top}$ fixed (adopted) & ${F_{gt}}$ & 0.323 & [0.309, 0.338] & 0.0072 \\
$\alpha_{top}$ fixed (adopted) & ${L_q}$ & 0.456 & [0.434, 0.478] & 0.011 \\
$\alpha_{top}$ fixed (adopted) & ${f_{htg}}$ & 1.83 & [1.617, 2.052] & 0.11 \\
$\alpha_{top}$ fixed (adopted) & $\chi^2$ / dof & 34.9 & 24 & reduced 1.46 \\
\bottomrule
\end{tabular}
\end{table}


The bed-side coefficient is validated instead on the plant data, where the furnace model supplies the outer-wall condition and
$f_{htg}$ is a free parameter. Fig.~\ref{fig:latham} and Tables~\ref{tab:latham}--\ref{tab:cal} summarise the calibration. The
adopted fit ($\alpha_{top}$ fixed) gives $F_{gt} = \calFgt{}$, $L_q = \calLq{}$ and $f_{htg} = \calFhtg{}$ with tight intervals, a
$\chi^2$ of \calChi{} on 24 degrees of freedom (reduced $\chi^2$ \calRedChi{}), and reproduces the outlet temperatures of all
three cases within 4 K, the lower wall temperatures within 1.1 K and the upper wall temperatures within 12.2 K. The four-parameter
fit lowers $\chi^2$ to 30.5 but is degenerate in $L_q$ and $\alpha_{top}$ (Table~\ref{tab:cal}); fixing $\alpha_{top}$ raises
$\chi^2$ by \SI{14}{\percent}, more than the \SI{10}{\percent} we set as the threshold for preferring the simpler model, so both
are reported and the three-parameter fit is adopted for all subsequent work because its parameters are identifiable
\todo{confirm the choice of the alpha-top-fixed configuration for the RQ work}. The heat-transfer factor of \calFhtg{} corresponds
to a bed-side coefficient of about \alphaIplant{} W m$^{-2}$ K$^{-1}$ at mid-height of the Plant A tube (four-parameter fit), i.e.\
between the two literature correlations and four times the value that reproduces the Xu--Froment figure; the range 300--3000
W m$^{-2}$ K$^{-1}$ is the epistemic span carried into the uncertainty analysis through $f_{htg}$.

Leave-one-case-out cross-validation (last two columns of Table~\ref{tab:latham}) shows that the held-out upper-wall temperature is
predicted within 5--18 K and the lower within 2 K. The upper-wall error is systematic: Plant C2, at \SI{71}{\percent} rate, is
predicted too hot at the upper elevation, and Plant A too cold. The upper peep hole lies in the flame zone, where a single
parabolic release profile with a load-independent flame length cannot follow the change in flame shape with firing rate. This is
the main structural limitation of the furnace model, and it is why the automated tests for the upper-wall error are recorded as
expected failures against the 4.5 K target derived from the stated measurement uncertainty. The consequence for the conclusions is
examined in Section~\ref{sec:rq4}.

\begin{figure}[htbp]
\centering
\includegraphics[width=0.6\textwidth]{fig06_creep_life_profile}
\caption{Outer-wall and mid-wall temperature profiles of the calibrated Plant A tube and the corresponding local rupture time from
the placeholder Larson--Miller curve at the local hoop stress; the hot spot \todo{state the hot-spot elevation} governs the life. Source:
\texttt{latham\_fit.yaml} (creep\_life\_first\_estimate), \texttt{yeh2021\_manaurite\_xm.yaml}, \rdtversion.}
\label{fig:creepprofile}
\end{figure}

Fig.~\ref{fig:creepprofile} shows the life profile of the calibrated Plant A tube. The hot spot lies at about \todo{quote z/L of the
hot spot from the figure or the YAML} of the heated length, where the outer wall reaches \baseTwo{} K at a hoop stress of
\SI{12.9}{MPa}; the rupture time from the placeholder minimum curve is $10^{8}$ h at this point. This absolute value is not
credible as a tube life and illustrates why the placeholder curve must be replaced: the manufacturer's curve carries no data below
about 20 MPa and the low hoop stress of a thick tube at 30 bar places the operating point in the extrapolated region
\todo{comment on the stress range of the placeholder curve and the extrapolation}. All results below are therefore expressed as
ratios of consumption rates, which depend on the slope of the master curve at the operating point rather than on its absolute
position.

\subsection{RQ1: what governs hot-spot temperature and life}
\label{sec:rq1}

\begin{figure}[htbp]
\centering
\includegraphics[width=\textwidth]{fig07_operating_maps}
\caption{Closed-loop operating maps in load and steam-to-carbon ratio at constant outlet temperature (\SI{1105.55}{K}): (a) hot-spot
temperature, (b) $\log_{10}$ of the life consumption relative to the base case and (c) net hydrogen per tube with contours of the
firing factor required. The star marks the calibrated Plant A base. Source: \texttt{data/lhs\_runs/grid\_sc\_load\_v2.csv.gz}, \rdtversion.}
\label{fig:maps}
\end{figure}

\begin{figure}[htbp]
\centering
\includegraphics[width=\textwidth]{fig08_sobol}
\caption{First-order ($S_1$) and total-effect ($S_T$) Sobol indices with \SI{95}{\percent} confidence intervals for (a) the hot-spot
temperature and (b) $\log_{10} t_r$ at the hot spot, from \nSaltelli{} physics evaluations of a Saltelli design over the seven operating
inputs. Source: \texttt{data/design/sobol\_v2.json}, \rdtversion.}
\label{fig:sobol}
\end{figure}

\begin{table}[htbp]
\centering
\caption{Sobol first-order ($S_1$) and total-effect ($S_T$) indices with 95\,\% confidence half-widths from a Saltelli design ($N=256$, 4096 physics runs) over the v2 operating space. Source: \texttt{data/design/sobol\_v2.json} (v2).}
\label{tab:sobol}
\small
\resizebox{\textwidth}{!}{%
\begin{tabular}{lrrrrrr}
\toprule
Input & $S_1$ ($T_{wo,max}$) & $S_T$ ($T_{wo,max}$) & $S_1$ ($\log_{10}t_r$) & $S_T$ ($\log_{10}t_r$) & $S_1$ (CH$_4$ slip) & $S_T$ (CH$_4$ slip) \\
\midrule
steam-to-carbon & 0.29 $\pm$ 0.09 & 0.28 $\pm$ 0.05 & 0.25 $\pm$ 0.08 & 0.25 $\pm$ 0.04 & 0.04 $\pm$ 0.03 & 0.04 $\pm$ 0.01 \\
load & 0.03 $\pm$ 0.03 & 0.03 $\pm$ 0.01 & 0.03 $\pm$ 0.03 & 0.03 $\pm$ 0.01 & 0.19 $\pm$ 0.07 & 0.19 $\pm$ 0.03 \\
firing & 0.34 $\pm$ 0.09 & 0.35 $\pm$ 0.06 & 0.29 $\pm$ 0.09 & 0.29 $\pm$ 0.05 & 0.54 $\pm$ 0.12 & 0.54 $\pm$ 0.10 \\
inlet $T$ & 0.06 $\pm$ 0.05 & 0.06 $\pm$ 0.01 & 0.05 $\pm$ 0.04 & 0.05 $\pm$ 0.01 & 0.08 $\pm$ 0.06 & 0.09 $\pm$ 0.02 \\
inlet $p$ & 0.02 $\pm$ 0.02 & 0.02 $\pm$ 0.00 & 0.16 $\pm$ 0.06 & 0.16 $\pm$ 0.03 & 0.03 $\pm$ 0.03 & 0.03 $\pm$ 0.01 \\
activity & 0.07 $\pm$ 0.04 & 0.08 $\pm$ 0.02 & 0.06 $\pm$ 0.05 & 0.06 $\pm$ 0.01 & 0.00 $\pm$ 0.01 & 0.00 $\pm$ 0.00 \\
excess air & 0.19 $\pm$ 0.06 & 0.19 $\pm$ 0.03 & 0.16 $\pm$ 0.07 & 0.16 $\pm$ 0.03 & 0.11 $\pm$ 0.04 & 0.11 $\pm$ 0.02 \\
\bottomrule
\end{tabular}}
\end{table}


The closed-loop maps (Fig.~\ref{fig:maps}) show the structure of the operating space when the outlet temperature is held by the
firing controller. At constant outlet temperature, the hot-spot temperature rises with load because the duty per tube rises while
the furnace-gas temperature at the flame zone is fixed by the fuel-to-feed ratio, and it falls with steam-to-carbon ratio because
the endothermic load moves towards the tube inlet \todo{check this physical reading of panel (a)}. The life map (panel b) is the
image of the temperature map through the Larson--Miller relation and spans about two decades across the space. Hydrogen production
(panel c) increases with both load and steam-to-carbon ratio, the latter through the shift of equilibrium, at the cost of a higher
firing factor.

The Sobol analysis over all seven inputs (Fig.~\ref{fig:sobol}, Table~\ref{tab:sobol}) ranks the drivers of the hot-spot temperature:
the specific firing factor ($S_T = \SobolFiringTwo{}$), the steam-to-carbon ratio (\SobolSCTwo{}) and the excess air (\SobolEAtwo{})
together explain \SI{82}{\percent} of the variance; inlet temperature, catalyst activity and load contribute below 0.1 each, and
inlet pressure is negligible. Excess air matters because at fixed fuel it lowers the adiabatic flame temperature and raises the
flue-gas flow, both of which move heat release down the tube. The ranking for $\log_{10} t_r$ is the same except that the inlet
pressure enters with $S_T = \SobolPlogtr{}$ through the hoop stress, which is invisible in the temperature. First-order and
total-effect indices agree within their confidence intervals for every input, so interactions are weak over the operating space and
the influence of each variable can be read from one-at-a-time maps. The methane slip is governed by the firing factor
($S_T = \SobolFiringSlip{}$) and the load, as expected for an approach-to-equilibrium quantity.

\begin{figure}[htbp]
\centering
\includegraphics[width=\textwidth]{fig09_surrogate_parity}
\caption{Parity of the selected surrogates against the physics model on the independent Saltelli test set for (a) the hot-spot
temperature, (b) dry methane slip and (c) $\log_{10} t_r$ at the hot spot, with nominal \SI{90}{\percent} CV+ conformal intervals.
Source: \texttt{data/design/surrogate\_metrics\_v2.json}, \rdtversion.}
\label{fig:parity}
\end{figure}

\begin{table}[htbp]
\centering
\caption{Best surrogate per target ($N_{train}$ = 2000 LHS, $N_{test}$ = 4096 Saltelli points), test $R^2$ and RMSE, 5-fold CV RMSE, and empirical coverage (PICP) and mean width (MPIW) of the nominal 90\,\% CV+ conformal interval. Source: \texttt{data/design/surrogate\_metrics\_v2.json}.}
\label{tab:surr}
\small
\resizebox{\textwidth}{!}{%
\begin{tabular}{lllrrrrr}
\toprule
Target & Unit & Model & $R^2$ & RMSE & CV RMSE & PICP$_{90}$ & MPIW$_{90}$ \\
\midrule
T\_wo\_max\_K & K & gp\_\_physics & 1 & 0.113 & 0.133 & 0.957 & 0.25 \\
z\_frac\_T\_wo\_max & -- & mlp\_\_base & 0.6436 & 0.0319 & 0.0522 & 0.971 & 0.0499 \\
T\_out\_K & K & gp\_\_physics & 1 & 0.0182 & 0.0254 & 0.940 & 0.0614 \\
P\_out\_bar & bar & gp\_\_base & 1 & 0.000287 & 0.000563 & 0.961 & 0.000998 \\
CH4\_slip\_dry\_pct & \% & gp\_\_physics & 1 & 0.00125 & 0.00142 & 0.953 & 0.00398 \\
H2\_net\_kmol\_h & kmol h$^{-1}$ & gp\_\_base & 1 & 0.000847 & 0.00116 & 0.946 & 0.00286 \\
duty\_W & W & gp\_\_base & 1 & 3.9e+03 & 4.14e+03 & 0.949 & 1.21e+04 \\
log10\_t\_r\_hot & -- & gp\_\_base & 0.9999 & 0.00566 & 0.00597 & 0.959 & 0.0111 \\
\bottomrule
\end{tabular}}
\end{table}


The surrogates that make the following campaigns possible are summarised in Table~\ref{tab:surr} and Fig.~\ref{fig:parity}. The
Gaussian process is selected for seven of the eight targets, and the physics features improve the hot-spot temperature, outlet
temperature and slip. Test RMSE is \surrRMSETwo{} K for the hot-spot temperature and \surrRMSElogtr{} decades for $\log_{10} t_r$,
two orders of magnitude below the calibration residuals, at a speed-up of \speedup{} relative to the physics model. The CV+
intervals over-cover: at a nominal level of 0.90 the empirical coverage is 0.94--0.97 for all targets (\PICPtwo{} for the hot-spot
temperature). Over-coverage is expected for CV+, whose guarantee is one-sided, and is conservative; the interval widths are
nevertheless small enough (\SI{0.25}{K} for the hot-spot temperature) that surrogate error is immaterial to the conclusions.

\subsection{RQ2: life cost of flexible operation and ageing}
\label{sec:rq2}

\begin{figure}[htbp]
\centering
\includegraphics[width=\textwidth]{fig11_scenarios_summary}
\caption{Annual life consumption per kmol of hydrogen relative to the steady base S1 for the RQ2 scenarios, from hourly accumulation
(bars) and from a single evaluation at the annual-mean inputs (diamonds). Source: \texttt{data/scenarios/summary\_v2.csv}, \rdtversion.}
\label{fig:scenarios}
\end{figure}

\begin{table}[htbp]
\centering
\caption{RQ2 scenario summary for one year (8760 h) of the Plant A twin at 21.6\,\% excess air: hydrogen per tube, Robinson damage $D$, life consumption per kmol H$_2$ relative to S1, ratio of the annual-mean-condition damage estimate to the hourly-accumulated damage, mean and maximum hot-spot temperature and mean firing factor. Life uses the Yeh Manaurite XM placeholder curve. Source: \texttt{data/scenarios/summary\_v2.csv}.}
\label{tab:scen}
\small
\resizebox{\textwidth}{!}{%
\begin{tabular}{lrrrrrrr}
\toprule
Scenario & H$_2$ (10$^3$ kmol) & $D$ (--) & Rel.\ life cost & Avg.-cond.\ ratio & $\bar T_{wo,max}$ (K) & max $T_{wo,max}$ (K) & Firing \\
\midrule
S1 steady base & 128.5 & 9.08e-05 & 1.000 & 1.000 & 1143 & 1143 & 0.978 \\
S2 daily load-following & 115.7 & 8.1e-05 & 0.991 & 0.982 & 1141 & 1143 & 0.944 \\
S3 renewable-following & 109.7 & 7.61e-05 & 0.983 & 0.982 & 1140 & 1146 & 0.927 \\
S4a mild over-firing (48 h) & 128.6 & 9.14e-05 & 1.007 & 0.993 & 1144 & 1157 & 0.979 \\
S4b severe over-firing (6 h) & 128.5 & 9.11e-05 & 1.004 & 0.996 & 1143 & 1179 & 0.978 \\
S4c hot band +40 K (24 h) & 128.5 & 9.33e-05 & 1.028 & 0.980 & 1144 & 1183 & 0.978 \\
S6 steam-to-carbon 2.5 & 119.9 & 9.83e-05 & 1.161 & 1.000 & 1144 & 1144 & 0.897 \\
S6 steam-to-carbon 3.5 & 139.8 & 8.23e-05 & 0.834 & 1.000 & 1142 & 1142 & 1.092 \\
S5 ageing, hold slip, year 1 & 128.5 & 6.49e-05 & 0.715 & 0.997 & 1138 & 1139 & 0.974 \\
S5 ageing, hold slip, year 4 & 128.5 & 0.000205 & 2.255 & 0.949 & 1156 & 1164 & 0.989 \\
S5 ageing, hold $T_{out}$, year 1 & 129.0 & 7e-05 & 0.768 & 0.998 & 1139 & 1140 & 0.982 \\
S5 ageing, hold $T_{out}$, year 4 & 127.4 & 0.000169 & 1.875 & 0.967 & 1153 & 1159 & 0.969 \\
\bottomrule
\end{tabular}}
\end{table}


Table~\ref{tab:scen} and Fig.~\ref{fig:scenarios} answer RQ2 for one year of operation at the calibrated base. Under the
outlet-temperature controller, load-following costs almost nothing in tube life per unit of hydrogen: the daily profile (S2) and the
renewable-following profile (S3) consume \SxDaily{} and \SxRenew{} of the S1 life per kmol H$_2$, i.e.\ slightly \emph{less}, because
the hot-spot temperature at constant outlet temperature falls with load (Fig.~\ref{fig:maps}a) and the hours at reduced load more
than compensate the hours at overload. Hydrogen output falls by \SI{10}{\percent} and \SI{15}{\percent} respectively, so the
absolute annual damage falls in proportion. A steady-state assessment at the annual-mean inputs reproduces the hourly accumulation
to within \SI{2}{\percent} for every scenario without discrete events (ratio \SxAvgErr{} for S3), which means that the non-linearity
of the Larson--Miller relation is not strong enough over the load range of this controller to make time-resolved simulation
necessary for smooth load variations \todo{discuss whether this conclusion holds for a plant without outlet-temperature control}.

Discrete over-firing events are different in kind. During a mild event (\SI{+10}{\percent} firing) the damage rate is \evMild{} times
the base rate, during a severe event (\SI{+25}{\percent}) it is \evSevere{} times, and a \SI{40}{K} hot band raises it by a factor of
\evHot{}; but because the events are short, their share of the annual damage is 0.4--2.7\,\% (\evHotShare{}\,\% for the hot band)
and the annual life cost per kmol rises by at most \SI{3}{\percent}. The annual-mean evaluation underestimates the event scenarios
(diamonds below bars in Fig.~\ref{fig:scenarios}), as it must for a convex damage law, but the error is small at these event
frequencies. Steam-to-carbon ratio has a larger steady effect: operating at 2.5 costs \SxSClow{} and at 3.5 costs \SxSChigh{} of the
S1 life per kmol H$_2$, because the higher ratio both cools the hot spot and raises the hydrogen yield.

Catalyst ageing dominates every operational effect. With the activity falling from 0.30 to 0.10 over four years and the controller
holding the methane slip, the life cost per kmol H$_2$ rises from \SxAgeSlipYone{} in the first year (fresh catalyst) to \SxAgeSlip{}
in the fourth as the hot spot climbs from \TwoAgeStart{} to \TwoAgeEnd{} K; holding the outlet temperature instead gives
\SxAgeTout{} in the fourth year because the slip is allowed to rise. Over the campaign the life consumed in year 4 is three times
that in year 1, and the question of when to change the catalyst is therefore also a tube-life question
\todo{comment on whether the activity range 0.30--0.10 is representative of a campaign}.

\subsection{RQ3: Pareto-optimal regimes and the steam credit}
\label{sec:rq3}

\begin{figure}[htbp]
\centering
\includegraphics[width=\textwidth]{fig12_pareto}
\caption{Physics-verified NSGA-II Pareto set in (a) firebox fuel per kmol H$_2$ against life consumption per kmol H$_2$, (b) net
hydrogen per tube against life consumption and (c) hydrogen against fuel; open circles are members infeasible under the physics
model, and the marked regimes are listed in Table~\ref{tab:regimes}. Source: \texttt{data/optimization/pareto\_v2.csv},
\texttt{regimes\_v2.csv}, \rdtversion.}
\label{fig:pareto}
\end{figure}

\begin{table}[htbp]
\centering
\caption{RQ3 representative regimes (physics-verified) from the NSGA-II Pareto set (200 $\times$ 300 generations; 199 of 200 members feasible under the physics model). Fuel = firebox fuel per kmol H$_2$; life = consumption rate per kmol H$_2$ relative to the base (Yeh placeholder curve). Source: \texttt{data/optimization/regimes\_v2.csv}.}
\label{tab:regimes}
\small
\resizebox{\textwidth}{!}{%
\begin{tabular}{lrrrrrrrrrr}
\toprule
Regime & Load & S/C & $T_{in}$ (K) & Excess air (\%) & Firing & H$_2$ (kmol h$^{-1}$) & Fuel (MJ kmol$^{-1}$) & Rel.\ life cost & $T_{wo,max}$ (K) & CH$_4$ slip (\%) \\
\midrule
base & 1.00 & 2.89 & 885 & 21.6 & 1.000 & 14.82 & 158.4 & 1.000 & 1147 & 7.81 \\
min life iso H2 & 0.96 & 4.00 & 895 & 17.5 & 0.873 & 14.72 & 133.7 & 0.083 & 1109 & 7.77 \\
knee & 1.10 & 4.00 & 900 & 5.0 & 0.871 & 17.22 & 130.5 & 0.280 & 1130 & 7.15 \\
max H2 & 1.10 & 4.00 & 900 & 5.0 & 1.150 & 19.26 & 154.1 & 4.329 & 1177 & 3.89 \\
min fuel overall & 0.60 & 4.00 & 900 & 5.0 & 0.850 & 9.87 & 121.3 & 0.472 & 1128 & 5.72 \\
min life overall & 0.60 & 4.00 & 858 & 24.9 & 0.850 & 9.18 & 130.4 & 0.046 & 1091 & 7.81 \\
\bottomrule
\end{tabular}}
\end{table}


The Pareto set (Fig.~\ref{fig:pareto}, Table~\ref{tab:regimes}) contains \nPareto{} members, of which \nFeasible{} remain feasible
under the physics model; the maximum discrepancy between surrogate and physics on the set is 0.03 in relative life and below
\SI{0.1}{MJ\,kmol^{-1}} in fuel. The calibrated base is dominated by \nDominating{} members. The knee of the front raises hydrogen
by \kneeHtwo{}\,\% while cutting fuel per kmol by \kneeFuel{}\,\% and life consumption per kmol to \kneeLife{} of the base; it does so
by running at \SI{110}{\percent} load with the minimum excess air, the maximum inlet temperature and a lower firing factor. At
constant production (iso-H$_2$ regime), the life cost can be reduced to \isoLife{} of the base with \isoFuel{}\,\% fuel by reducing
the firing and letting the hot spot fall from \baseTwo{} to \isoTwo{} K. The maximum-hydrogen corner (\maxHtwoGain{}\,\%) costs
\maxHtwoLife{} times the base life per kmol and runs close to the wall-temperature limit.

Every member of the front sits at the upper bound of the steam-to-carbon ratio (3.96--4.00). This is an artefact of the objective:
with firebox fuel alone, additional steam is free, and the equilibrium shift and the cooling of the hot spot both favour maximum
steam. The steam-credit sweep (Fig.~\ref{fig:steam}, Table~\ref{tab:steam}) shows how the answer changes when steam raising and feed
preheat are charged. With steam fully charged ($\alpha = 0$) the minimum-fuel regime moves to a steam-to-carbon ratio of
\steamMinFuelZeroSC{} and the knee to 3.3, while the minimum-life regime stays at 4.0 but now costs \steamMinLifeZeroFuel{}\,\% more
total fuel than the base for a life cost of \steamMinLifeZero{}. For $\alpha \ge 0.5$, i.e.\ when at least half of the steam duty is
recovered from the flue gas or credited as export steam, the ratio of 4.0 is optimal for all three regimes. The plant's own steam
balance therefore decides whether tube life and fuel economy pull in the same direction or against each other
\todo{relate the steam-credit range to a typical SMR steam balance with waste-heat boiler and export steam}.

\begin{figure}[htbp]
\centering
\includegraphics[width=\textwidth]{fig13_steam_credit}
\caption{Iso-production front as a function of the steam credit $\alpha$: (a) optimal steam-to-carbon ratio and (b) life
consumption per kmol H$_2$ relative to the base for the minimum-life, knee and minimum-fuel regimes. Source:
\texttt{data/optimization/steam\_credit\_v2.json}, \rdtversion.}
\label{fig:steam}
\end{figure}

\begin{table}[htbp]
\centering
\caption{Steam-credit sweep of the iso-production (H$_2$ = base $\pm$1\,\%) life--fuel front: $\alpha=1$ counts firebox fuel only, $\alpha=0$ charges steam raising and feed preheat in full. Physics-verified. Source: \texttt{data/optimization/steam\_credit\_v2.json}.}
\label{tab:steam}
\small
\resizebox{\textwidth}{!}{%
\begin{tabular}{llrrrrrrr}
\toprule
$\alpha$ & Regime & S/C & Load & Firing & Excess air (\%) & Rel.\ life cost & Total fuel vs base (\%) & $T_{wo,max}$ (K) \\
\midrule
0.00 & min life & 3.99 & 0.96 & 0.909 & 24.6 & 0.059 & +3.2 & 1103 \\
0.00 & knee & 3.29 & 0.97 & 0.850 & 5.1 & 0.828 & -6.6 & 1144 \\
0.00 & min fuel & 2.66 & 1.00 & 0.888 & 5.0 & 4.539 & -8.6 & 1171 \\
0.25 & min life & 4.00 & 0.96 & 0.943 & 24.9 & 0.057 & +2.3 & 1103 \\
0.25 & knee & 3.80 & 0.94 & 0.850 & 5.0 & 0.365 & -7.1 & 1131 \\
0.25 & min fuel & 2.92 & 0.99 & 0.865 & 5.0 & 2.016 & -8.9 & 1158 \\
0.50 & min life & 4.00 & 0.96 & 1.039 & 25.0 & 0.057 & +5.5 & 1103 \\
0.50 & knee & 4.00 & 0.95 & 0.850 & 13.8 & 0.103 & -7.6 & 1112 \\
0.50 & min fuel & 3.63 & 0.95 & 0.850 & 5.0 & 0.480 & -10.5 & 1135 \\
0.75 & min life & 4.00 & 0.96 & 0.911 & 25.0 & 0.057 & -6.6 & 1103 \\
0.75 & knee & 4.00 & 0.96 & 0.850 & 13.5 & 0.105 & -12.2 & 1112 \\
0.75 & min fuel & 4.00 & 0.93 & 0.850 & 5.0 & 0.276 & -14.6 & 1127 \\
1.00 & min life & 4.00 & 0.96 & 0.958 & 25.0 & 0.057 & -7.3 & 1103 \\
1.00 & knee & 4.00 & 0.96 & 0.850 & 15.3 & 0.087 & -17.7 & 1109 \\
1.00 & min fuel & 4.00 & 0.93 & 0.850 & 5.0 & 0.276 & -20.4 & 1127 \\
\bottomrule
\end{tabular}}
\end{table}


\subsection{RQ4: uncertainty and robustness}
\label{sec:rq4}

\begin{figure}[htbp]
\centering
\includegraphics[width=\textwidth]{fig14_uq_distributions}
\caption{Monte Carlo distributions (\nMC{} samples of all uncertainty groups through the physics model) of (a) the hot-spot
temperature, with the \SI{1193}{K} limit, and (b) $\log_{10} t_r$ at the hot spot, at the S1 base, the RQ3 knee, the iso-H$_2$
minimum-life regime and the end of the fourth ageing year. Source: \texttt{data/uq/mc\_v2\_*.csv.gz}, \rdtversion.}
\label{fig:uq}
\end{figure}

\begin{table}[htbp]
\centering
\caption{RQ4 Monte Carlo (scrambled Sobol design, $N=2048$, all uncertainty groups) through the physics model: median and 5--95\,\% interval of the hot-spot temperature, of $\log_{10}t_r$ and of the paired life consumption per kmol H$_2$ relative to the base. Source: \texttt{data/uq/uq\_summary\_v2.json}.}
\label{tab:uq}
\small
\resizebox{\textwidth}{!}{%
\begin{tabular}{llll}
\toprule
Operating point & $T_{wo,max}$ (K) & $\log_{10}t_r$ (h) & Rel.\ life cost \\
\midrule
S1 base & 1145 [1127, 1169] & 7.91 [7.04, 8.73] & 1.000 [1.000, 1.000] \\
RQ3 knee & 1132 [1107, 1165] & 8.30 [7.33, 9.20] & 0.338 [0.201, 0.643] \\
RQ3 min-life iso-H$_2$ & 1110 [1088, 1138] & 8.90 [7.94, 9.80] & 0.100 [0.069, 0.159] \\
S5 year 4, hold slip & 1165 [1138, 1203] & 7.35 [6.28, 8.30] & 3.478 [1.841, 8.974] \\
S/C 2.5 & 1146 [1129, 1166] & 7.88 [7.04, 8.69] & 1.153 [0.916, 1.350] \\
S/C 3.5 & 1144 [1127, 1172] & 7.94 [7.04, 8.83] & 0.846 [0.587, 1.119] \\
load 0.70 & 1138 [1123, 1156] & 8.08 [7.22, 8.90] & 0.966 [0.638, 1.350] \\
load 0.85 & 1141 [1124, 1162] & 8.00 [7.14, 8.82] & 0.951 [0.787, 1.143] \\
\bottomrule
\end{tabular}}
\end{table}

\begin{table}[htbp]
\centering
\caption{Variance decomposition at the S1 base point: one-group-at-a-time share of the joint variance ($N=1024$ per group) and group total-effect Sobol index from a Gaussian-process surrogate of the joint sample. Source: \texttt{data/uq/uq\_summary\_v2.json}.}
\label{tab:var}
\small
\resizebox{\textwidth}{!}{%
\begin{tabular}{lrrrr}
\toprule
Group & share $\log_{10}t_r$ & $S_T$ $\log_{10}t_r$ & share $T_{wo,max}$ & $S_T$ $T_{wo,max}$ \\
\midrule
kinetics & 0.233 & 0.225 & 0.513 & 0.620 \\
heat transfer & 0.050 & 0.114 & 0.111 & 0.302 \\
creep & 0.622 & 0.655 & 0.211 & 0.093 \\
measurement & 0.094 & 0.036 & 0.202 & 0.092 \\
sum / interaction residual & 0.999 / +0.001 &  & 1.037 / -0.037 &  \\
\bottomrule
\end{tabular}}
\end{table}

\begin{table}[htbp]
\centering
\caption{Robustness of the RQ2/RQ3 conclusions: fraction of paired Monte Carlo samples supporting each statement, with the calibrated flame length $L_q$ and with the structural variant $L_q\,\mathrm{load}^{0.5}$. Source: \texttt{data/uq/uq\_summary\_v2.json}.}
\label{tab:robust}
\small
\begin{tabular}{lrr}
\toprule
Statement & nominal $L_q$ & $L_q \propto \mathrm{load}^{0.5}$ \\
\midrule
ageing year 4 costs $>2\times$ S1 life per kmol & 0.906 & 0.906 \\
knee regime $<0.5\times$ base & 0.806 & 0.936 \\
S/C 3.5 cheaper in life than S/C 2.5 & 0.836 & 0.836 \\
S2 daily $\le$ S1 & 0.568 & 0.000 \\
S3 renewable $\le$ S1 & 0.632 & 0.001 \\
\bottomrule
\end{tabular}
\end{table}


Table~\ref{tab:uq} and Fig.~\ref{fig:uq} give the propagated uncertainty. At the base point the \SI{90}{\percent} interval of the
hot-spot temperature is \uqTwoWidth{} K wide and that of $\log_{10} t_r$ spans \uqLogtrLo{}--\uqLogtrHi{}, i.e.\ a factor of 50 in
absolute rupture time. The paired ratios are much tighter: the knee regime costs 0.20--0.64 of the base life per kmol and the
fourth ageing year 1.8--9.0 times, so the ranking of regimes survives an uncertainty that makes the absolute life meaningless. The
nominal result lies inside the \SI{90}{\percent} interval at every point. The variance decomposition (Table~\ref{tab:var},
Fig.~\ref{fig:variance}) attributes \shareCreep{}\,\% of the variance of $\log_{10} t_r$ to the creep group, dominated by the
heat-to-heat scatter and the Larson--Miller constant, 23\,\% to the kinetics, 9\,\% to the wall-temperature measurement and only
\shareHT{}\,\% to the calibrated heat-transfer parameters; for the hot-spot temperature the kinetics group dominates
(\shareKinTwo{}\,\%), followed by the measurement offset and the wall thickness. The one-at-a-time shares sum to within 0.04 of one,
and the Sobol total effects on the surrogate confirm the ranking, so the groups act almost additively.

\begin{figure}[htbp]
\centering
\includegraphics[width=0.55\textwidth]{fig15_variance_shares}
\caption{One-group-at-a-time shares of the variance of $\log_{10} t_r$ and of the hot-spot temperature at the S1 base point.
Source: \texttt{data/uq/uq\_summary\_v2.json}, \rdtversion.}
\label{fig:variance}
\end{figure}

The robustness fractions (Table~\ref{tab:robust}) state which conclusions hold across the joint uncertainty. That the fourth ageing
year costs more than twice the S1 life per kmol holds in \robAge{} of the samples, that the knee regime halves the life cost in
\robKnee{}, and that steam-to-carbon 3.5 is cheaper in life than 2.5 in \robSC{}; these three conclusions are robust and, for the knee,
become more so under the flame-length variant. The conclusion that flexible operation costs no life per kmol is not robust: it holds
in \robFlexNom{} of the samples with the calibrated flame length, close to a coin toss because the nominal effect is within
\SI{2}{\percent} of unity, and it is reversed in all samples when the flame shortens with load, under which the hot spot at
\SI{70}{\percent} load moves up the tube into the flame zone and the life cost per kmol rises to \lqvarLoadLow{} times the base. The
sign of the flexible-operation effect is therefore determined by a structural feature of the furnace, the response of the flame
length to load, that the three plant cases cannot identify.
